\ifdefined\HMTTRITFormal
\chapter{État local ternaire et orientation du transport}
\label{chap:trit}
\index{TRIT}

Le TRIT est la coordonnée ternaire qui conserve l'orientation locale d'une
transition. Sa définition précède toute interprétation temporelle : les
trois valeurs distinguent des types algébriques de transport et n'acquièrent une
lecture de passé, de seuil ou de futur que lorsqu'un paramètre ordonné est fixé.

\section{Signature ternaire orientée et espace d'états}

L'espace des signatures visibles est
\[
\TRIT=\{-1,0,+1\}.
\]
Ses éléments sont des symboles orientés, non des résultats d'une opération dans
\(\ZZ\). L'application équilibrée
\[
\FF_3\longrightarrow\TRIT,
\qquad0\mapsto0,\quad1\mapsto+1,\quad2\mapsto-1,
\]
permet de passer du résidu ternaire à la signature orientée.

\begin{definition}[État complet du TRIT et fibre de retenue]
\label{def:trit-levantado-carry-rev8}
Dans un pas local dont l'amplitude requiert au plus une retenue, la
décomposition ternaire équilibrée s'écrit
\[
 \widehat\tau=(r,c),\qquad
 a+b=r+3c,
 \qquad r,c\in\{-1,0,+1\}.
\]
La première coordonnée \(r\) est le résidu visible et la seconde \(c\) conserve
le quotient qui a rendu cette lecture possible.  En particulier, la fibre au-dessus du
zéro visible contient trois états distincts :
\[
 0^+=(0,+1),\qquad 0^0=(0,0),\qquad 0^-=(0,-1),
 \qquad
 \operatorname{vis}(0^+)=\operatorname{vis}(0^0)
 =\operatorname{vis}(0^-)=0.
\]
\end{definition}

Cette séparation locale entre résidu et quotient est celle qui, dans l'élévation
de Bockstein--Hensel ultérieure, prend la forme matricielle
\[
 x_k A=u_k+3x_{k+1}.
\]
Ici \(u_k\) publie le résidu du niveau \(k\), tandis que \(x_{k+1}\)
transporte la mémoire nécessaire pour récupérer le relèvement \(x_k\).  Ainsi,
un zéro visible ne détermine à lui seul ni l'état ni son futur admissible.
Cette affirmation est algébrique : elle ne transforme pas \(0^\pm\) en un infinitésimal,
une énergie, une torsion physique ou un intervalle temporel.

Un trit visible ne détermine pas l'ensemble des continuations admissibles.
L'espace complet des états de transport contient les champs
\[
x=(t,\mu,o,h,c,\sigma,\lambda,g).
\]
Ici \(t\in\TRIT\) est la signature visible ;
\(\mu\in\mathsf M:=\{\Sigma,\Pi,\varnothing\}\), l'étiquette de l'observable
additive, multiplicative ou neutre ; \(o\), l'orientation ; \(h\), la coordonnée d'extension ;
\(c\), la retenue accumulée ; \(\sigma\), la signature de frontière ; \(\lambda\), le
registre de trajectoire ; et
\(g\in\ZZ/9\ZZ\), la phase. La projection \(x\mapsto t\) élimine toutes les
autres composantes.

\begin{remark}[Distinction entre symbole et état]
Deux états peuvent afficher le même trit et posséder des continuations distinctes.
Un mot ternaire ne détermine pas non plus à lui seul l'étiquette d'observable,
l'orientation, la retenue ou la
monodromie. Cette distinction explique pourquoi cinq régions de \(\pi\) peuvent
partager un mot \(w_\pi\) sans devenir une seule région.
\end{remark}

\fi

\ifdefined\HMTTRITDevelopments
\chapter{Algèbres quadratiques et régimes de transport}
\label{chap:trit-algebras}

\section{\(3\) algèbres quadratiques associées}

Pour \(\tau\in\TRIT\), soit \(\bar\tau\in\{-1,0,+1\}\subset\RR\) son
image par l'identification équilibrée précédente. Nous considérons l'algèbre réelle
unitaire
\[
\mathbb A_\tau=\RR[X]/(X^2+\bar\tau),
\]
dans laquelle \(J\) désigne la classe de \(X\) et \(I\) l'unité. Cette convention
sépare le symbole TRIT de sa valeur scalaire.
Selon le signe, on obtient trois types :
\[
\begin{array}{ccl}
\tau=+1&:&J^2=-I,\quad\text{type elliptique ou complexe},\\
\tau=0&:&J^2=0,\quad\text{type parabolique ou nilpotent},\\
\tau=-1&:&J^2=+I,\quad\text{type hyperbolique ou scindé}.
\end{array}
\]
Les exponentielles montrent la différence :
\[
e^{tJ}=
\begin{cases}
\cos t+J\sin t,&J^2=-I,\\
1+tJ,&J^2=0,\\
\cosh t+J\sinh t,&J^2=+I.
\end{cases}
\]
Les trois algèbres sont, respectivement, l'extension complexe, l'algèbre des
nombres duaux et l'algèbre réelle déployée. Leurs sous-groupes exponentiels
produisent des évolutions elliptiques, paraboliques et hyperboliques sans introduire de
géométrie continue supplémentaire dans la définition du TRIT\@.
Les réalisations réelles concrètes sélectionnées par le Coxeter d'APP, le
générateur nilpotent et l'incidence des feuillets TPK sont construites, sans
dupliquer cette classification abstraite, dans
\cref{sec:tres-generadores-concretos-trit}.

\begin{figure}[H]
\centering
\resizebox{\textwidth}{!}{\begin{tikzpicture}[font=\small,box/.style={draw,rounded corners=2mm,minimum width=4.2cm,minimum height=4.4cm,align=center,inner sep=3mm}]
  \node[box,draw=hmtviolet,fill=hmtlightviolet] (ell) at (0,0) {};
  \node[font=\bfseries,align=center,text width=3.6cm,hmtviolet]
    at (0,1.55) {Elliptique /\\complexe};
  \draw[hmtviolet,thick,->] (0,-.25) arc (-90:255:.62);
  \node[fill=hmtlightviolet,inner sep=1.5pt] at (0,.32) {\(J^2=-I\)};
  \node[align=center] at (0,-1.18) {sous-groupe compact\\retour elliptique};

  \node[box,draw=hmtgray,fill=hmtpaper] (par) at (5,0) {};
  \node[font=\bfseries,align=center,text width=3.6cm,hmtgray]
    at (5,1.55) {Parabolique /\\nilpotente};
  \draw[hmtgray,thick,->] (4.1,.3) .. controls (4.6,-.2) and (5.4,-.2) .. (5.9,.3);
  \draw[hmtgray,dashed] (4.05,.3)--(5.95,.3);
  \node at (5,.72) {\(J^2=0\)};
  \node[align=center] at (5,-1.18) {flot unipotent\\état seuil};

  \node[box,draw=hmtgold,fill=hmtlightgold] (hyp) at (10,0) {};
  \node[font=\bfseries,align=center,text width=3.6cm,hmtgold]
    at (10,1.55) {Hyperbolique /\\dédoublée};
  \draw[hmtgold,thick,->] (9.35,-.55) .. controls (9.62,-.10) and (9.62,.35) .. (9.78,.60);
  \draw[hmtgold,thick,->] (10.65,-.55) .. controls (10.38,-.10) and (10.38,.35) .. (10.22,.60);
  \node[fill=hmtlightgold,inner sep=1.5pt] at (10,.05) {\(J^2=+I\)};
  \node[align=center] at (10,-1.18) {deux sous-espaces propres\\évolution hyperbolique};

  \node[font=\footnotesize,align=center,hmtgray] at (5,-2.75) {L'interprétation temporelle est obtenue en fixant une paramétrisation ordonnée des trois types algébriques.};
\end{tikzpicture}
}
\caption{Trois types algébriques associés à l'état ternaire orienté. La
première ligne rassemble les identités algébriques ; la seconde, leur interprétation
par rapport à une paramétrisation ordonnée.}
\label{fig:trit-atlas}
\end{figure}

\section{Compatibilité et trajectoires paramétrées}

L'expression \emph{All Possible Paths} invite à imaginer des chemins parcourus dans
le temps. HMT situe l'objet une étape en amont. Soit \(\Omega_{\rm HMT}\) un
ensemble d'états locaux muni d'une relation de transition admissible
\(\leadsto\), d'une signature TRIT et de données de transport qui seront typées dans la
section suivante. Si
\(\mathcal T\) est un ensemble discrètement ordonné, une trajectoire
paramétrée est une application
\[
\gamma:\mathcal T\longrightarrow\Omega_{\rm HMT}
\]
qui satisfait \(\gamma(t)\leadsto\gamma(t^+)\) pour chaque successeur \(t^+\).
L'ordre de \(\mathcal T\) transforme la compatibilité en une suite
paramétrée. Avant de choisir un ordre, il n'y a qu'une famille compatible
indexée par le graphe ou l'ensemble partiellement ordonné correspondant ; on ne l'appelle section que
si une projection \(p:E\to\mathcal T\) a été déclarée avec
\(p\circ\gamma=\operatorname{id}_{\mathcal T}\).
Par conséquent, une somme finie de chemins peut se lire comme une somme sur les
temporalisations d'une structure qui, à son niveau premier, ne présuppose pas
un temps extérieur.

\section{Paramétrisation ordonnée de la compatibilité}

Une paramétrisation ordonnée permet d'interpréter le retour, le seuil et l'ouverture
comme trois positions relatives dans une évolution. La convention utilisée dans
\cref{fig:trit-atlas} est
\[
\begin{array}{ccl}
J^2=-I&\rightsquigarrow&\text{retour et orientation inverse},\\
J^2=0&\rightsquigarrow&\text{état seuil},\\
J^2=+I&\rightsquigarrow&\text{décomposition en deux directions propres}.
\end{array}
\]
Dans la paramétrisation temporelle adoptée, \(+1\) représente
le retour, la mémoire et le passé ; \(0\), le seuil et le présent ; et \(-1\), l'ouverture
bidirectionnelle et le futur. Cette lecture ordonnée ne modifie pas les trois types
quadratiques précédents.
L'involution exacte, typée pour chaque \(n\geq1\),
\[
C_{\rm rev}:\TRIT^n\longrightarrow\TRIT^n,
\qquad
C_{\rm rev}(\tau_1,\ldots,\tau_n)=(-\tau_n,\ldots,-\tau_1)
\]
échange \(+1\) et \(-1\), préserve \(0\) et satisfait
\(C_{\rm rev}^2=\operatorname{id}_{\TRIT^n}\).
Ainsi, l'orientation temporelle n'est pas introduite comme un quatrième état : elle résulte
de la paramétrisation de la signature avec une involution qui conserve le
seuil.

\section{Décomposition spectrale du type hyperbolique}

Dans le type hyperbolique, les projecteurs
\[
P_\pm=\frac{I\pm J}{2}
\]
séparent deux directions propres. L'involution qui change l'orientation les
échange. Algébriquement, \(P_+P_-=0\), \(P_++P_-=I\) et
\(JP_\pm=\pm P_\pm\). Le type hyperbolique contient donc deux sous-espaces
propres complémentaires et une opération exacte qui les permute.

\begin{remark}[Structure complète de l'état TRIT]
TRIT ne remplace pas les signes \(-,0,+\). C'est l'observable minimale d'un
état avec étiquette d'observable, orientation et registre de trajectoire.
Ses trois algèbres associées permettent
qu'un même support admette des fermetures elliptiques, des seuils paraboliques et des ouvertures
hyperboliques avant de fixer une évaluation temporelle.
\end{remark}
\fi
